% TOP_PROBLEM: {"id": 2896, "title": "Kirby Problem 4.20", "category_id": 7, "category": {"id": 7, "name": "topology", "display_name": "Topology", "description": "Properties preserved under continuous deformations.", "slug": "topology", "order_index": 7, "created_at": "2026-07-31T15:26:25.671Z"}, "status": "open"}
% FIRST_POSTED: null
% ENTRY_KIND: statement_only
% Imported from: batch13.tex; UnsolvedMath ID 2896
\documentclass[11pt]{article}
\usepackage[margin=25mm]{geometry}
\usepackage{fontspec}
\setmainfont{Latin Modern Roman}
\usepackage{amsmath,amssymb,mathtools}
\usepackage{unicode-math}
\setmathfont{Latin Modern Math}
\usepackage{xurl,hyperref}
\hypersetup{colorlinks=true,urlcolor=blue}
\setcounter{secnumdepth}{0}
\setlength{\parindent}{0pt}
\setlength{\parskip}{5pt}
\setlength{\emergencystretch}{3em}
\newcommand{\sref}[2]{\hyperlink{source:#1:#2}{[S#2]}}
\newcommand{\cataloguescope}[1]{\paragraph{Recorded target.}#1}
\begin{document}
% UnsolvedMath ID 2896; source catalogue code KP-4.20
\setcounter{section}{2895}
\section{Kirby Problem 4.20}\label{2896}
{\small\sffamily UnsolvedMath ID 2896 · Topology}
\subsection{Definitions and mathematical statement}

\paragraph{Shared notation.}$\mathbb N=\{1,2,\ldots\}$, and $\mathbb Z,\mathbb Q,\mathbb R,\mathbb C$ denote the integers, rationals, reals and complex numbers. Any other conventions are specified locally.


For a closed topological four-manifold $M$, write $\operatorname{ks}(M)\in\mathbb Z/2\mathbb Z$ for the Kirby--Siebenmann obstruction evaluated on its mod-two fundamental class; it obstructs a compatible piecewise-linear structure and therefore a smooth structure. Smoothable means homeomorphic to a smooth manifold.

Let $\star\mathbb{RP}^4$ denote the unique homeomorphism class of topological manifolds homotopy equivalent to $\mathbb{RP}^4$ with $\operatorname{ks}=1$. This is Ruberman's nonsmoothable model. Its identification is given in K3, Problem 4.87, Remark (1), and Problem 4.88, Remark (2), PDF page 261. The star symbol here names that model; it is not a claim that every star construction satisfies a connected-sum cancellation identity.

Let $\mathrm{En}$ be the complex Enriques surface, obtained as the quotient of a K3 surface by a free holomorphic involution. A K3 surface is a simply connected compact complex surface with trivial canonical bundle. Let $\star\mathrm{En}$ be its oriented topological star partner: it has the oriented homotopy type of $\mathrm{En}$ and opposite Kirby--Siebenmann invariant, hence $\operatorname{ks}=1$. Existence of this named partner is part of the original K3 statement; uniqueness within that oriented homotopy type follows from Teichner's classification theorem for fundamental group $\mathbb Z/2\mathbb Z$. \sref{2896}{3}
\paragraph{Statement recorded in the supplied source.}
Are the topological connected sums
\[
(\star\mathbb{RP}^4)\#(\star\mathbb{RP}^4),\qquad
(\star\mathrm{En})\#(\star\mathrm{En})
\]
smoothable? The source also records the same question for the mixed sum
\[
(\star\mathbb{RP}^4)\#(\star\mathrm{En}).
\]
Each sum has zero Kirby--Siebenmann invariant by additivity, but the requested conclusion is smoothability without adding $S^2\times S^2$ factors. \sref{2896}{2}
\subsection{Short English statement}
Are the doubles of the nonsmooth projective-space model and the Enriques star partner smoothable, and is their mixed connected sum smoothable?
\subsection{Sources}
\begin{description}
\item[\hypertarget{source:2896:1}{S1}] ULAM.AI and the UnsolvedMath Contributors, UnsolvedMath: A Curated Collection of Open Mathematics Problems, record 2896 (KP-4.20). CC BY 4.0. \url{https://huggingface.co/datasets/ulamai/UnsolvedMath}.
\item[\hypertarget{source:2896:2}{S2}] R. İnanç Baykur, Robion C. Kirby and Daniel Ruberman, K3—A New Problem List in Low-Dimensional Topology, Mathematical Surveys and Monographs 295 (2026), author version, Problem 4.20 and Remarks, PDF page 205; also Problems 4.87--4.88, PDF page 261. \url{https://math.berkeley.edu/sites/default/files/surv-295-ruberman-watermarked-author-pdf.pdf}.
\item[\hypertarget{source:2896:3}{S3}] Peter Teichner, On the star-construction for topological 4-manifolds, Introduction and Theorem 1, pages 1--2. \url{https://math.berkeley.edu/~teichner/Papers/star.pdf}.
\end{description}
\label{end:2896}
\end{document}
